\section{Descent of the actual invariant Frobenius matrices}
\label{sec:invariant-frobenius-descent-v055}

\status{This chapter proves the full original inhomogeneous matrix
descent theorem, on the supplied Wach coefficient ring and the actual
Witt tensor ring. It also proves the abstract additive and matrix
bijectivity statements used in that descent. The approximation argument
uses the exact sharp-root coordinates of an invariant tilt element and
the previously proved completed-fixed-field theorem. Together with WIS18,
this leaves the actual Wach prism as the one remaining period-ring
supplier in the conditional all-layer return WLD26. The original tied
period data and crystalline-realization input are retained.}

There are two different approximations in this proof. Powers of the
cyclotomic parameter control the zeroth Witt coordinate; powers of $p$
control the successive Witt coordinates. Frobenius improves the first
approximation. Saturation then permits division by $p$, so that the
argument can be repeated. The final passage to the coefficient ring uses
its $p$-adic closed image, already proved in WIS19.

Keep the original prime $p$, period data $P,AI$, finite fields $K,LL$,
level $s$, and supplied records $\mathrm{Good}(K,s,LL)$,
$\Gamma,R,I,\mathrm{IoW}(I)$. In the notation of WIS put
\[
 O=\mathcal O_{LL},\quad S=\mathcal O_C^\flat,\quad F=\operatorname{Frac}(S),
 \quad B=O\otimes_{\mathbb Z_p}W(S),\quad W=O\otimes_{\mathbb Z_p}W(F).
\]
Write $c:B\hookrightarrow W$, $A=\prod_{i\in\mathbb Z/n\mathbb Z}O[[T]]$,
$E=R.\mathrm{EE}$, $a:A\to E$, $i:A\to B$ and $j=I.\mathrm{ioE}:E\to W$.
Thus $ja=ci$. The symbol $T$ in $A$ means the tuple whose every entry
is the power-series variable. Put
\[
 \tau_s=\epsilon_s-1\in S,\quad \mu_s=[\epsilon_s]-1\in W(S),\quad
 t=c(1\otimes\mu_s)=j(a(T)),\qquad
 \mathcal N_m=t^mc(B)+pW.                                      \tag{IFD.1}
\]
Here $t$ is a unit of $W$. Accordingly $\mathcal N_m$ must not be
replaced by the ideal $t^mW+pW$. The coefficient $c(B)$ in (IFD.1)
is indispensable. The fixed ring used at the endpoint is
\[
 W^H,\qquad H=G_K\cap\ker\chi.
                                                               \tag{IFD.2}
\]
The finite coefficient basis in WIS4 is kept throughout. Its coordinates
will be denoted by $w=\sum_\alpha b_\alpha\otimes w_\alpha$.
The formulas and completeness statements of WIS18--WIS25 are already
human proofs and are used at their stated scopes.

\subsection{The filtration and its closed coefficient image}

% BEGIN IFD1
\paragraph{IFD1: the elementary near-parameter calculus.}
\label{claim:IFD1}
More generally let $c:B_0\to W_0$ be a ring map, let
$t=c(\tau)$, and put $N_m=t^mc(B_0)+pW_0$ for $m\geq0$.
Then $N_m$ is an additive subgroup, $N_{m+1}\subseteq N_m$,
$N_mN_k\subseteq N_{m+k}$, and $c(B_0)N_m\subseteq N_m$.
If $t$ is a unit, then
\[
             t^k x\in N_{m+k}\quad\Longrightarrow\quad x\in N_m.
                                                               \tag{IFD.3}
\]
For the actual rings, membership in $\mathcal N_m$ is equivalent to
$(w_\alpha)_0\in\tau_s^m S$ for every $\alpha$, with $S$ embedded in $F$.

\begin{proof}
Addition and negation act on the two coefficients in a presentation
$x=t^mc(b)+pr$. Descending the exponent multiplies $b$ by $\tau$.
Multiplying two such presentations gives leading term
$t^{m+k}c(bb')$ and a remainder divisible by $p$; the assertion about
$c(B_0)$ is the same calculation. In (IFD.3), divide the presentation
$t^kx=t^{m+k}c(b)+pr$ by the unit $t^k$.
The last assertion is exactly WIS23: in the unchanged finite $O$-basis,
zeroth Witt coordinates detect the ideal $pW$, and the zeroth coordinate
of $\mu_s$ is $\tau_s$. This proves the equivalence with the precise
subgroup in (IFD.1).
\end{proof}
% END IFD1

% BEGIN IFD2
\paragraph{IFD2: Frobenius improves the filtration and detects constants.}
\label{claim:IFD2}
The original Frobenius $\phi$ satisfies
\[
 \phi(c(B))\subseteq c(B),\qquad
 \phi(\mathcal N_m)\subseteq\mathcal N_{pm},\qquad
 \phi(t)\in\mathcal N_2.                                      \tag{IFD.4}
\]
In particular it preserves every $\mathcal N_m$. If the image in $W$
of $o\in O$ belongs to $\mathcal N_1$, then $o\in pO$.
These are all four clauses of the original
\texttt{frobenius\_and\_constants} theorem.

\begin{proof}
Frobenius on the Witt factor is induced by $x\mapsto x^p$ on $S$ or
$F$, and it fixes the $O$ factor. It therefore preserves $c(B)$.
If $(w_\alpha)_0=\tau_s^m v_\alpha$, then
$(\phi w_\alpha)_0=\tau_s^{mp}v_\alpha^p$.
IFD1 proves the middle inclusion. Since $t\in\mathcal N_1$ and
$p\geq2$, it proves the third assertion and also preservation of
each subgroup. These arguments include $p=2$ and every $s\geq0$.

Expand $o=\sum_\alpha z_\alpha b_\alpha$ with $z_\alpha\in\mathbb Z_p$.
Its $\alpha$th Witt coordinate is the image of $z_\alpha$ in $W(F)$.
If $p\nmid z_\alpha$, then $z_\alpha$ is a unit of $\mathbb Z_p$,
so its image in $W(S)$, and then its zeroth coordinate in $S$, are
units. That coordinate cannot lie in $\tau_s S$: WIS22 implies
that $\tau_s$ is a nonunit, since otherwise $1$ would lie in every
$\tau_s^mS$. IFD1 thus forces $p\mid z_\alpha$ for every $\alpha$.
Dividing the finitely many coefficients gives $o\in pO$.
\end{proof}
% END IFD2

% BEGIN IFD3
\paragraph{IFD3: a near-divisible component idempotent is zero.}
\label{claim:IFD3}
Let $c:B_0\hookrightarrow W_0$ be injective, with $p$-saturated image,
where $W_0$ has no $p$-torsion and is $p$-adically separated.
Let $t=c(\tau)$ be a unit, and suppose $\bigcap_mN_m\subseteq pW_0$.
If $e,c_0,d,w\in B_0$ and $N\geq0$ satisfy
\[
 e^2=e,\qquad p=c_0d,\qquad d^{N+1}=pw,
                 \qquad c(c_0e)\in N_1,
                                                               \tag{IFD.5}
\]
then $e=0$.

\begin{proof}
Write $c(c_0e)=tc(b)+pr$. Saturation gives $r=c(r_0)$ and
injectivity gives $c_0e=\tau b+pr_0$ in $B_0$.
Multiplication by $d$ yields
$p(e-dr_0)=\tau db$. In $W_0$ this says
$c(db)=p t^{-1}c(e-dr_0)$. Saturation gives an $s\in B_0$ with
$c(s)=t^{-1}c(e-dr_0)$, and hence
\[
                              e=dr_0+\tau s.
\]
Raise this equality to the power $N+1$. Idempotence gives
$e^{N+1}=e$, the term $(dr_0)^{N+1}$ is divisible by $p$, and every
other term is divisible by $\tau$. Thus $e\in\tau B_0+pB_0$.
Raising this last expression to any positive power shows
$c(e)\in N_m$ for every $m$. The intersection hypothesis gives
$c(e)\in pW_0$, and saturation and injectivity give $e=pe_1$ in $B_0$.
For every $k$, idempotence now gives
$e=e^{k+1}=p^ke_1^ke$. Consequently $c(e)\in p^kW_0$ for all $k$.
Separatedness and then injectivity give $e=0$.
\end{proof}
% END IFD3

% BEGIN IFD4
\paragraph{IFD4: detection of every component's constant coefficient.}
\label{claim:IFD4}
In addition to IFD3, let $O_0$ be a discrete valuation domain with
$p\ne0$, let $B_0$ be an $O_0$-algebra, and let
$i_0:\prod_{i\in\mathbb Z/n\mathbb Z}O_0[[T]]\hookrightarrow B_0$
be an injective $O_0$-algebra map.
If a tuple of constants $u=(u_i)$ satisfies
$c(i_0(u))\in N_1$, then $p\mid u_i$ for every $i$.
This is the full generic component-constant implication
\texttt{Rational.const\_dvd}.

\begin{proof}
Suppose $p\nmid u_i$. Total divisibility in the valuation domain gives
$p=u_i d$, and $d$ is a nonunit: if it were a unit, $p$ would divide
$u_i$. If $p$ has valuation $e_p$, some power $d^{N+1}$ is divisible
by $p$. Let $\varepsilon_i$ be the nonzero idempotent which is $1$
in component $i$ and $0$ elsewhere. Multiplication of the assumed
near-relation by $i_0(\varepsilon_i)$ gives
$c(u_i i_0(\varepsilon_i))\in N_1$. Apply IFD3 with
$e=i_0(\varepsilon_i)$ and with the images of $u_i,d$ and
$d^{N+1}/p$. It gives $i_0(\varepsilon_i)=0$, contrary to injectivity.
The argument uses the actual component idempotent and does not require
the tensor ring $B_0$ to be a domain.
\end{proof}
% END IFD4

% BEGIN IFD5
\paragraph{IFD5: the precise integral-series normal form.}
\label{claim:IFD5}
Let $A_0=\prod_iO_0[[T]]$, with ring maps $i_0:A_0\to B_0$ and
$c:B_0\to W_0$. Suppose $t=c(i_0(T))$ is a unit and suppose the
constant-coefficient conclusion of IFD4 holds, without requiring its
particular proof hypotheses. Then
\[
 c(i_0(f))\in N_m\quad\Longrightarrow\quad
                   f=T^mf_1+pf_2\quad(f_1,f_2\in A_0).        \tag{IFD.6}
\]
For any sequence $b_k\in A_0$ there is also a single $\sigma\in A_0$
such that
\[
             \sigma-\sum_{k<m}T^kb_k\in T^mA_0
                                  \quad\hbox{for every }m.  \tag{IFD.7}
\]

\begin{proof}
For (IFD.6), induct on $m$. At $m=0$ take $f_1=f,f_2=0$.
At the next step the inductive expression is $f=T^mf_1+pf_2$.
Use IFD1 to cancel $t^m$ from the stronger assumption in $N_{m+1}$;
it gives $c(i_0(f_1))\in N_1$. Write $f_1=Tg+u$, where $u$ is
its tuple of constant coefficients. Since $c(i_0(Tg))\in N_1$,
so is $c(i_0(u))$. The assumed criterion gives $u=pu'$.
Thus $f=T^{m+1}g+p(f_2+T^mu')$, proving the induction.

For (IFD.7), define the coefficient of $T^r$ in component $i$ of
$\sigma$ to be the coefficient in $\sum_{k\leq r}T^kb_k(i)$.
This is a finite sum. For $r<m$ that coefficient equals the one in
$\sum_{k<m}T^kb_k(i)$, since the additional terms have degree at
least $m$. Their difference therefore has its first $m$ coefficients
zero, which is exactly divisibility by $T^m$ in every component.
No analytic convergence is used in this formal-series sum.
\end{proof}
% END IFD5

% BEGIN IFD6
\paragraph{IFD6: full near-parameter closedness modulo the prime.}
\label{claim:IFD6}
Keep the hypotheses of IFD5. Let $a_0:A_0\to E_0$ and
$j_0:E_0\to W_0$ be ring maps with $j_0a_0=ci_0$.
Assume that for every $e\in E_0$ there are $f\in A_0$, $k\geq0$
and $r\in E_0$ with
\[
                     e a_0(T)^k=a_0(f)+pr,                 \tag{IFD.8}
\]
and that $\bigcap_mN_m\subseteq pW_0$. Then
\[
 \bigl(\forall m\ \exists e_m\in E_0:\ x-j_0(e_m)\in N_m\bigr)
          \quad\Longrightarrow\quad x\in j_0(E_0)+pW_0.    \tag{IFD.9}
\]
This is the full original generic \texttt{near\_closed} statement.

\begin{proof}
First, if $j_0(e)\in N_m$, (IFD.8) implies
$c(i_0(f))\in N_{m+k}$. IFD5 gives
$f=T^{m+k}f_1+pf_2$. After applying the maps and dividing by the
unit $t^k$, we get the useful normal form
\[
                     j_0(e)=c(i_0(T^mf_1))+pr'.             \tag{IFD.10}
\]
Choose the approximants $e_m$ in (IFD.9). Their differences satisfy
$j_0(e_{m+1}-e_m)\in N_m$. Apply (IFD.10) to write
\[
                     j_0(e_{m+1}-e_m)=c(i_0(T^mb_m))+pr_m.
\]
Choose the single formal sum $\sigma$ from (IFD.7).
For every $M$, telescoping the displayed differences shows that
$x-j_0(e_0)-c(i_0(\sigma))$ is the sum of $x-j_0(e_M)$,
a term in $t^Mc(B_0)$, and a term in $pW_0$.
It therefore belongs to every $N_M$ and hence to $pW_0$.
Since $c(i_0(\sigma))=j_0(a_0(\sigma))$, this proves (IFD.9).
In particular this step has not assumed completeness of $E_0$
for the parameter $T$.
\end{proof}
% END IFD6

% BEGIN IFD7
\paragraph{IFD7: the actual image is fixed, has bounded poles, and is near-closed.}
\label{claim:IFD7}
For the original maps fixed above,
\[
 j(E)\subseteq W^H,\qquad \phi(j(E))\subseteq j(E),\qquad
 \forall e\in E\ \exists k:\ t^kj(e)\in\mathcal N_0.       \tag{IFD.11}
\]
Moreover
\[
 \forall m\ \exists e_m:\ x-j(e_m)\in\mathcal N_m
               \quad\Longrightarrow\quad x\in j(E)+pW.    \tag{IFD.12}
\]
All the original \texttt{ioE\_fixed}, \texttt{bounded\_pole} and
\texttt{closed} conclusions hold.

\begin{proof}
For $h\in H$, the specified character identity and the injection
$O\hookrightarrow LL$ give $\chi_\Gamma(\operatorname{res}_\Gamma h)=1$.
Its stipulated injectivity gives $\operatorname{res}_\Gamma h=1$.
The original $\Gamma$-equivariance of $j$ therefore says $hj(e)=j(e)$.
Its Frobenius identity proves the second assertion.
The record $R.\mathrm{EE\_dense}$, at precision $p$, gives (IFD.8)
for the original $a:A\to E$. Applying $j$ proves the bounded-pole
assertion.

The map $i:A\to B$ is injective: if $i(f)=i(g)$, the identity
$ja=ci$ and WIS18 give $a(f)=a(g)$, and the supplied injection $a$
gives $f=g$. WIS5 gives injectivity and $p$-saturation of $c$,
and absence of $p$-torsion and $p$-adic separation in $W$.
The image of $T$ is a unit because $a(T)$ is a unit in the original
record. WIS24 gives $\bigcap_m\mathcal N_m=pW$.
Finally WIS1 makes $O$ a DVR with $p\ne0$.
Thus IFD4 supplies the component-constant criterion for these exact
maps, and IFD6 proves (IFD.12). The same reasoning proves the generic
\texttt{closed} implication with its original printed hypotheses:
each actual WIS input is simply replaced by the corresponding supplied
injectivity, saturation, torsion, separation and intersection premise.
\end{proof}
% END IFD7

\subsection{Invariant tilt approximation by the original generators}

For this subsection $K$ is any finite extension, with no
$\mathrm{Good}$ assumption. Let $K_0$ be its maximal unramified
subfield, $K_r=K_0(\zeta_{p^r})$ and
$K_\infty=\bigcup_rK_r$. Write $H_\infty=G_{K_\infty}$.
The roots here are the original compatible roots in $AI$.
Let $G_0\subseteq S$ be the subring generated by all $\epsilon_r$
and by the prime-to-$p$ root tilts in the source:
\[
 G_0=\mathbb Z[\,\epsilon_r\ (r\geq0),\ x_b:
  \operatorname{IsTiltOf}(u,m,x_b),\ u\in K_0,\ u^m=1,
                  \ m>0,\ p\nmid m\,].                    \tag{IFD.13}
\]
As a subring of a characteristic-$p$ ring, the indicated integer
scalars of course factor through $\mathbb F_p$.
No closure for a topology is part of this definition.

% BEGIN IFD8
\paragraph{IFD8: sharp-root coordinates preserve the exact fixed group.}
\label{claim:IFD8}
For $x\in S$ write $x^{\sharp,N}\in\mathcal O_C$ for its exact
$N$th sharp-root coordinate. For every original Galois element $g$,
\[
        (g x)^{\sharp,N}=g(x^{\sharp,N}).                    \tag{IFD.14}
\]
If $x$ is fixed by $H_\infty$, then each $x^{\sharp,N}$ belongs
to the completion $\widehat{K_\infty}$ and has norm at most one.
If the $N$th inverse-limit coordinate of $z\in S$ is zero, then
\[
                            |z|_\flat\leq |p|^{p^N}.        \tag{IFD.15}
\]

\begin{proof}
Choose arbitrary integral representatives $a_r$ of the inverse-limit
coordinates $x_r\in\mathcal O_C/p$. The coordinate formula is
\[
                 x^{\sharp,N}=\lim_{r\to\infty}a_{N+r}^{p^r}.
                                                               \tag{IFD.16}
\]
Here is the estimate which makes both the limit and its independence
of representatives precise. If $|u|,|v|\leq1$ and $|u-v|\leq|p|^k$
with $k\geq1$, the binomial expansion gives
$|u^p-v^p|\leq|p|^{k+1}$. The intermediate binomial coefficients
are divisible by $p$, and the last term has valuation at least
$pk\geq k+1$. Iterating gives
$|u^{p^r}-v^{p^r}|\leq|p|^{k+r}$.
Since $a_{r+1}^p\equiv a_r\pmod p$, successive terms of (IFD.16)
differ by at most $|p|^{r+1}$. They are Cauchy; changing any
representatives has the same vanishing bound. This includes $p=2$.

The original field \texttt{actT\_coeff} says that $g(a_r)$ is a
representative of $(gx)_r$. The original \texttt{actC\_norm} says
that $g$ is an isometry of $C$, hence continuous. Apply it to
(IFD.16) and pass through the limit to obtain (IFD.14).
This uses only the original action fields, not an additional choice
of action on a completion.

If $x$ is $H_\infty$-fixed, (IFD.14) fixes each sharp-root
coordinate. LS1, applied to this same isometric extension of the
action on $\Omega$, gives
$C^{H_\infty}=\widehat{\Omega^{H_\infty}}
=\widehat{K_\infty}$. The last equality is the algebraic
fixed-field correspondence for $H_\infty=G_{K_\infty}$.
The coordinates in (IFD.16) are integral, so the limit has norm at
most one. Finally $z^{\sharp,N}$ reduces to the $N$th coordinate
of $z$. If it is zero modulo $p$, its norm is at most $|p|$.
The exact relation $(z^{\sharp,N})^{p^N}=z^\sharp$ proves (IFD.15).
No additivity of the sharp map has been asserted or used.
\end{proof}
% END IFD8

% BEGIN IFD9
\paragraph{IFD9: integral coefficients in the finite cyclotomic layers.}
\label{claim:IFD9}
The nonzero norms in $K_0$ are integral powers of $p$.
Every integral element of $K_0$ is congruent modulo $p$ to zero or
to a root of unity in $K_0$ of order prime to $p$.
For every $r\geq1$,
\[
                 \mathcal O_{K_r}=\mathcal O_{K_0}[\zeta_{p^r}].
                                                               \tag{IFD.17}
\]

\begin{proof}
We recall the elementary unramified facts in the form needed here.
For a finite residue field $\mathbb F_{p^f}$, lift a monic
irreducible polynomial for one residue generator to $\mathbb Z_p[X]$.
Its unramified integral extension is finite free over $\mathbb Z_p$;
the residue is $\mathbb F_{p^f}$, and its maximal ideal is generated
by $p$. To see the last assertion and the norm statement directly,
expand an element in the lifted basis. The smallest common power
of $p$ in its coefficients leaves a nonzero residue. Such a residue
is a unit: lift its inverse and sum the geometric series for the
error in $p$. Completeness holds coordinatewise. Thus this ring is
a DVR with uniformizer $p$.

Every nonzero residue has a unique lift satisfying $u^{p^f-1}=1$.
For existence, Newton correction for $X^{p^f-1}-1$ is integral
and doubles the error valuation, because its derivative is a unit
at the chosen residue. The corrections converge by completeness.
For uniqueness, if two such roots have the same residue, division
of their power difference by their difference gives a unit
(its residue is the nonzero derivative), so the difference is zero.
These roots generate the unramified extension. Conversely any
prime-to-$p$ root of unity lies in one of these extensions: choose
$f$ with its order dividing $p^f-1$ and use the unique simple-root
lifts of all the roots in the finite residue field. The union of
these fields is the original maximal unramified extension.
Its finite subfield $K_0$ is therefore unramified, with the stated
norm group and residue representatives. Equivalently, finite
subfields correspond to the subfields of the finite residue fields;
the lifted residue basis gives the same degree and the same
uniformizer $p$.

Let $e=p^{r-1}(p-1)$ and $\pi=\zeta_{p^r}-1$.
The elementary cyclotomic product calculation in LS2 gives
$|\pi|=p^{-1/e}$: primitive conjugates have the same distance from
$1$ and their product at $1$ is $p$. For $a_i\in K_0$ the nonzero
terms $a_i\pi^i$, $0\leq i<e$, have distinct valuations modulo
$\mathbb Z$. Hence they cannot cancel. In particular these $e$
powers are linearly independent over $K_0$. The cyclotomic
polynomial gives the opposite degree bound, so they are a basis
of $K_r$ over $K_0$.

Write an integral $z\in K_r$ in this basis. The norm of its sum is
the largest of the distinct term norms, so
$|a_i|p^{-i/e}\leq1$ for every $i$. Since the valuation of $a_i$
is an integer and $0\leq i/e<1$, each $a_i$ is integral.
Thus $z\in\mathcal O_{K_0}[\pi]=\mathcal O_{K_0}[\zeta_{p^r}]$.
The reverse inclusion follows from integrality of the generator,
proving (IFD.17), including $p=2,r=1$ where $e=1$.
\end{proof}
% END IFD9

% BEGIN IFD10
\paragraph{IFD10: every completed integral element has a generator lift.}
\label{claim:IFD10}
For every $N\geq0$ and $z\in\widehat{K_\infty}$ with $|z|\leq1$,
there exists $w\in G_0$ such that
\[
                  w_N=z\bmod p\quad\hbox{in }\mathcal O_C/p.
                                                               \tag{IFD.18}
\]
The same coefficient index $N$ and the same unramified field $K_0$
are used throughout.

\begin{proof}
Let $U_N$ be the inverse image in $\mathcal O_C$ of the image of
$G_0$ under its $N$th coefficient map. It is a subring containing
$p\mathcal O_C$. It is therefore closed under replacing an element
by any other integral element congruent to it modulo $p$.

The $N$th coefficient of $\epsilon_j$ is
$\zeta_{p^{j+N}}\bmod p$. Raising to $p^N$ shows that
$\zeta_{p^j}\in U_N$ for every $j$.
If $u\in K_0$, $u^m=1$, $p\nmid m$, choose an integer
$a$ with $ap\equiv1\pmod m$. The compatible sequence
\[
                           u^{p^N a^k}\qquad(k\geq0)
\]
defines a tilt of $u^{p^N}$ of order dividing $m$ and has $N$th
coordinate $u$. It is one of the generators allowed in (IFD.13).
Thus $u\in U_N$. IFD9 and congruence stability now imply
$\mathcal O_{K_0}\subseteq U_N$ and then
$\mathcal O_{K_j}\subseteq U_N$ for every $j\geq1$.

Choose $z'\in K_\infty$ with $|z-z'|<|p|$. The ultrametric
inequality and $|z|\leq1$ give $|z'|\leq1$. It lies in some finite
layer, so $z'\in U_N$. Congruence stability gives $z\in U_N$.
This is exactly (IFD.18). The construction requires only membership
in the subring generated by the specified elements; it does not
silently replace that subring by its completion.
\end{proof}
% END IFD10

% BEGIN IFD11
\paragraph{IFD11: the full original invariant tilt approximation.}
\label{claim:IFD11}
If $x\in S$ is fixed by $H_\infty$, then for every $m\geq0$
there exists $w\in G_0$ such that
\[
                           (\epsilon-1)^m\mid x-w
                                      \quad\hbox{in }S.    \tag{IFD.19}
\]
This proves the original \texttt{TiltInvariants.invariant\_approx}.

\begin{proof}
The element $\epsilon-1$ is nonzero. Indeed its sharp-root
coordinates reduce to $\zeta_{p^r}-1$, whose norm is greater than
$|p|$ for some $r$ by IFD9; their reductions cannot all be zero.
Set $b=(\epsilon-1)^m$, which is nonzero also for $m=0$.
Choose $N$ so large that $|p|^{p^N}\leq|b|_\flat$.
By IFD8 the exact root $x^{\sharp,N}$ belongs to
$\widehat{K_\infty}\cap\mathcal O_C$.
IFD10 gives $w\in G_0$ whose $N$th coefficient equals its
reduction modulo $p$, hence equals the $N$th coefficient of $x$.
The $N$th coefficient of the tilt difference $x-w$ is zero.
IFD8 therefore gives $|x-w|_\flat\leq|p|^{p^N}\leq|b|_\flat$.
The tilt valuation-ring property proved in LS4 says exactly that
this norm inequality is divisibility by $b$ in $S$.
This proves (IFD.19). The sharp map is used only for norms,
limits and multiplicative root relations, never as an additive map.
\end{proof}
% END IFD11

% BEGIN IFD12
\paragraph{IFD12: fractional invariant approximation with a common pole.}
\label{claim:IFD12}
Let $y\in F$ be fixed by $H_\infty$, and keep any $s,m\geq0$.
There is an $N_0$ such that for every $N\geq N_0$ one can find
$u\in G_0$, $v\in S$ with
\[
                         \tau_s^N y=u+\tau_s^{N+m}v.
                                                               \tag{IFD.20}
\]
This is the complete original \texttt{flat\_invariant\_approx} scope.

\begin{proof}
The element $\tau_s$ is nonzero, has norm less than one, and is
fixed by $H_\infty$: its inverse-limit coordinates are the
cyclotomic roots, all fixed by that group, so the coordinate action
fixes it. The norm and nonvanishing follow from IFD9 and the same
coordinate test as IFD11. Thus for all sufficiently large $N$,
$\tau_s^N y$ is integral; it is also fixed.
Apply IFD11 to it with exponent $N+m$.
In the characteristic-$p$ ring $S$,
$\epsilon-1=(\epsilon_s-1)^{p^s}=\tau_s^{p^s}$.
Consequently divisibility by $(\epsilon-1)^{N+m}$ implies
divisibility by $\tau_s^{N+m}$, giving (IFD.20).
The threshold $N_0$ depends only on $y$ and $s$, and every larger
$N$ is allowed. A single maximum therefore works for any finite
list of fractional coordinates, as required below.
\end{proof}
% END IFD12

% BEGIN IFD13
\paragraph{IFD13: the correct fixed coordinates in a good layer.}
\label{claim:IFD13}
Under the original $\mathrm{Good}(K,s,LL)$, if $y\in W^H$ then
every $(y_\alpha)_0\in F$ is fixed by $H_\infty$.

\begin{proof}
The original good-layer condition includes $K\subseteq K_s$.
Every element of $H_\infty$ fixes $K_s$, hence $K$, and fixes
all cyclotomic roots, hence has cyclotomic character $1$.
Therefore $H_\infty\subseteq H$.
The action on the $O$ factor is the identity. In the fixed finite
basis, applying an element of $H_\infty$ to the equality $hy=y$
gives equality of each Witt coordinate, and then of its zeroth
coordinate. These are the actual actions on $W(F)$ and $F$.
No replacement of $K$ by an unramified extension has occurred.
\end{proof}
% END IFD13

% BEGIN IFD14
\paragraph{IFD14: elements eventually in the coefficient image form a subring.}
\label{claim:IFD14}
Let $j_0:E_0\to W_0$ be a ring map and $\phi_0:W_0\to W_0$
a ring endomorphism with $\phi_0(j_0(E_0))\subseteq j_0(E_0)$.
Then
\[
 \mathcal P=\{x\in W_0:\ \exists r\geq0,\
                    \phi_0^r(x)\in j_0(E_0)+pW_0\}         \tag{IFD.21}
\]
is a subring. If its defining condition holds for $r$, it holds
for every larger exponent.

\begin{proof}
A ring endomorphism fixes the integer $p$. Applying $\phi_0$
to $j_0(e)+pw$ again gives an element of $j_0(E_0)+pW_0$,
which proves persistence. Zero, one and negatives satisfy (IFD.21).
For two elements satisfying it at exponents $r,r'$, pass to
$\max(r,r')$. Their images then lie in the subring
$j_0(E_0)+pW_0$, which is closed under addition and multiplication.
The sum and product therefore satisfy the same condition.
This verifies the exact original \texttt{perf} subring without
assuming that $j_0(E_0)$ is already perfect.
\end{proof}
% END IFD14

% BEGIN IFD15
\paragraph{IFD15: approximation inside any subring containing the generators.}
\label{claim:IFD15}
Let $T_0\subseteq W$ be a subring containing the image of $O$,
$t^{-1}$, and the elements $1\otimes[x]$ for every generator $x$
in (IFD.13). Under the original good-layer condition, for each
$y\in W^H$ and $m\geq0$ there is $y_0\in T_0$ with
$y-y_0\in\mathcal N_m$.

\begin{proof}
Consider the inverse image of $T_0$ in $W(S)$ under
$f_W:w\mapsto1\otimes w$, followed by its image under the
zeroth-coordinate ring map to $S$. This image is a subring of $S$
containing every generator of $G_0$. Thus every $u\in G_0$ has
a lift $w\in W(S)$ with $w_0=u$ and $f_W(w)\in T_0$.
This step does not claim that Teichmuller lifting is additive.

By IFD13 the finitely many coordinates $(y_\alpha)_0$ are
$H_\infty$-fixed. Choose a common $N$ above all the thresholds
in IFD12. It gives $u_\alpha\in G_0,v_\alpha\in S$ with
\[
       \tau_s^N(y_\alpha)_0
                      =u_\alpha+\tau_s^{N+m}v_\alpha.
\]
Choose lifts $w_\alpha$ in the preceding inverse image and put
$z=\sum_\alpha b_\alpha\otimes w_\alpha$ in $W$.
Each term is the product of the image of $b_\alpha\in O$ and
$f_W(w_\alpha)$, so $z\in T_0$. IFD1's exact coordinate
criterion gives $t^Ny-z\in\mathcal N_{N+m}$.
Set $y_0=t^{-N}z\in T_0$ and use (IFD.3) to obtain
$y-y_0\in\mathcal N_m$. This proves the original
\texttt{approx\_subring} implication for every supplied $T_0$.
\end{proof}
% END IFD15

% BEGIN IFD16
\paragraph{IFD16: the actual invariant approximation has eventual coefficient values.}
\label{claim:IFD16}
For every $y\in W^H$ and $m\geq0$ there are $r\geq0$,
$y_0\in W$, $e\in E$ and $v\in W$ with
\[
          y-y_0\in\mathcal N_m,\qquad \phi^r(y_0)=j(e)+pv.
                                                               \tag{IFD.22}
\]
All the original \texttt{Rational.approx} hypotheses and outputs
are retained.

\begin{proof}
Use the subring $\mathcal P$ of IFD14 for the actual $j,\phi$.
It contains $j(E)$, hence the image of $O$ and $t^{-1}$.
If $x=\epsilon_r$, then the Witt Frobenius identity gives
\[
 \phi^r(1\otimes[\epsilon_r])
       =1\otimes[\epsilon]
       =(1+t)^{p^s}=j\bigl((1+a(T))^{p^s}\bigr).
\]
Thus these generators belong to $\mathcal P$.
For a prime-to-$p$ root tilt $x_b$ of $u\in K_0$, each
$\Gamma.\mathrm{emb}_i(u)$ is integral in $LL$, since its positive
power is $1$. The original \texttt{IoW} identity identifies
$1\otimes[x_b]$ with the image under $i$ of the tuple of these
integral constants. Applying $c$ identifies it with an element of
$j(E)$. It too belongs to $\mathcal P$.
IFD15 now approximates $y$ in $\mathcal P$ at the prescribed
precision. Membership in $\mathcal P$ supplies $r,e,v$ in
(IFD.22). There is no assertion that $y_0$ itself is fixed; the
later argument requires only the fixedness of $y$ and the displayed
approximation.
\end{proof}
% END IFD16

\subsection{An additive contraction, followed by prime-adic lifting}

% BEGIN IFD17
\paragraph{IFD17: one prime-adic step for the kernel.}
\label{claim:IFD17}
We state the exact general setting used in IFD17--IFD20.
Let $N$ be an abelian group, $p$ a natural number,
$\Phi:N\to N$ additive, $D\subseteq M\subseteq N$ subgroups,
and $F_m\subseteq N$ a decreasing sequence of subgroups.
Assume the following properties:
\begin{enumerate}
\item $\Phi(M)\subseteq M$, $\Phi(D)\subseteq D$ and
      $\Phi(F_m)\subseteq F_{m+1}$.
\item $M$ and $D$ are $p$-saturated in $N$: if $px$ belongs
      to either subgroup, so does $x$.
\item If $f_i\in M\cap F_i$ for all $i$, there is $s\in M$
      with $s-\sum_{i<m}f_i\in F_m$ for every $m$.
\item If $x$ can be approximated by $D$ modulo every $F_m$,
      then $x\in D+pN$.
\item For each $y\in M$, there are $j\geq0$ and $d\in D$
      with $\Phi^jy-d\in F_0$.
\item Every sequence $x_i\in M$ has a sum $s\in M$ satisfying
      $s-\sum_{i<k}p^ix_i\in p^kN$ for all $k$.
\item If $x\in D+p^kN$ for every $k$, then $x\in D$.
\end{enumerate}
If $z\in M$ and $\Phi z-z\in D$, then there are $d\in D$,
$r\in M$ with
\[
                         z-d=pr,\qquad \Phi r-r\in D.      \tag{IFD.23}
\]

\begin{proof}
Put $\delta=\Phi z-z\in D$.
For every $j$, telescoping gives
$\Phi^jz-z=\sum_{i<j}\Phi^i\delta\in D$.
Property 5 therefore provides an approximation $z-d_0\in F_0$.
If $z-d_m\in F_m$, put $d_{m+1}=\Phi d_m-\delta\in D$.
Then $z-d_{m+1}=\Phi(z-d_m)\in F_{m+1}$.
Thus $z$ is approximable by $D$ at every filtration level.
Property 4 gives $z-d=pr$ for some $d\in D,r\in N$.
Since $pr\in M$, property 2 gives $r\in M$.
Finally
\[
                p(\Phi r-r)=(\Phi z-z)-(\Phi d-d)\in D.
\]
The saturation of $D$ gives the second part of (IFD.23).
Only the displayed hypotheses have been used; no topology on the
abstract group has been silently imposed.
\end{proof}
% END IFD17

% BEGIN IFD18
\paragraph{IFD18: injectivity on the quotient by the coefficient subgroup.}
\label{claim:IFD18}
Under the hypotheses of IFD17,
\[
                 z\in M,\quad \Phi z-z\in D
                                    \quad\Longrightarrow\quad z\in D.
                                                               \tag{IFD.24}
\]

\begin{proof}
Apply IFD17 repeatedly. Inductively one obtains
$z-d_k=p^kr_k$ with $d_k\in D$, $r_k\in M$ and
$\Phi r_k-r_k\in D$. At $k=0$ take $d_0=0,r_0=z$.
If $r_k-e_k=pr_{k+1}$ is the next step, set
$d_{k+1}=d_k+p^ke_k$. Then
$z-d_{k+1}=p^{k+1}r_{k+1}$ with the same required properties.
Property 7 of IFD17 now gives $z\in D$.
Equivalently, the map induced by $\Phi-1$ on $M/D$ is injective.
\end{proof}
% END IFD18

% BEGIN IFD19
\paragraph{IFD19: one prime-adic step for the cokernel.}
\label{claim:IFD19}
Under the same hypotheses, for every $y\in M$ there exist
$z\in M$, $d\in D$ and $r\in M$ such that
\[
                         y-(\Phi z-z)-d=pr.                \tag{IFD.25}
\]

\begin{proof}
Choose $j,d_0$ with $e=\Phi^jy-d_0\in F_0$ using property 5.
It also belongs to $M$. The terms $\Phi^ie$ belong to $M\cap F_i$.
Property 3 supplies $s\in M$ with
$s-S_m\in F_m$, where $S_m=\sum_{i<m}\Phi^ie$.
The finite telescoping identity
\[
                          \Phi S_m-S_m=\Phi^me-e
\]
gives
\[
 \Phi s-s+e=(\Phi(s-S_m)-(s-S_m))+\Phi^me\in F_m.
\]
Here $\Phi(F_m)\subseteq F_{m+1}\subseteq F_m$ has been used.
Property 4 applied with the zero approximant in $D$ gives
$\Phi s-s+e=d_1+pr$ for some $d_1\in D,r\in N$.
Set
\[
                     z=-\left(\sum_{i<j}\Phi^iy+s\right)\in M.
\]
A second telescoping calculation gives
$y-(\Phi z-z)=d_0+(\Phi s-s+e)=d_0+d_1+pr$.
This proves (IFD.25), with $d=d_0+d_1$.
Its left side belongs to $M$, so saturation puts $r$ in $M$.
The construction is the ordinary geometric-series argument, with
the sign determined by solving for $\Phi-1$.
\end{proof}
% END IFD19

% BEGIN IFD20
\paragraph{IFD20: the full original additive quotient bijectivity theorem.}
\label{claim:IFD20}
Under all the hypotheses printed in IFD17,
\[
                         \Phi-1:M/D\longrightarrow M/D
                                      \quad\hbox{is bijective}.       \tag{IFD.26}
\]
Explicitly, for every $y\in M$ there is $z\in M$ with
$y-(\Phi z-z)\in D$, and the kernel assertion is (IFD.24).
This is the full original
\texttt{sub\_one\_bijective\_quotient\_of\_topNilpotent}
at its original additive-group scope.

\begin{proof}
Starting with $y_0=y$, repeatedly apply IFD19 to obtain
\[
                  y_i-(\Phi z_i-z_i)-d_i=p y_{i+1},
                   \qquad z_i,y_{i+1}\in M,\ d_i\in D.
\]
For every $k$ the finite sum satisfies
\[
 y-\sum_{i<k}p^i(\Phi z_i-z_i)-\sum_{i<k}p^id_i=p^ky_k.
\]
Property 6 gives $z\in M$ with
$z-\sum_{i<k}p^iz_i=p^kr_k$ for each $k$.
Since $\Phi$ is additive, it commutes with integer multiplication.
Subtracting $(\Phi-1)$ of this last equality from the preceding
one yields
\[
 y-(\Phi z-z)-\sum_{i<k}p^id_i
                         =p^k\bigl(y_k-(\Phi r_k-r_k)\bigr).
\]
The sum of the $d_i$ terms lies in $D$. Property 7 therefore gives
$y-(\Phi z-z)\in D$. This proves surjectivity; IFD18 proves
injectivity. Both statements descend to the quotient because
$\Phi(D)\subseteq D$.
\end{proof}
% END IFD20

\subsection{The original matrix theorem and its actual specialization}

% BEGIN IFD21
\paragraph{IFD21: a finite coefficient matrix has one contracting filtration.}
\label{claim:IFD21}
Let $c:B_0\to W_0$, $t=c(\tau)$ and $N_m=t^mc(B_0)+pW_0$
be as in IFD1. Let $\phi_0$ preserve $c(B_0)$ and satisfy
$\phi_0(t)\in N_2$. Let $j_0:E_0\to W_0$ have the property
that each $j_0(e)$ has a bounded pole:
$t^kj_0(e)\in N_0$ for some $k$.
For every finite square matrix $P\in M_d(E_0)$ there is a
single $c_0\geq0$ such that the additive operator
\[
                   \Phi(v)=j_0(P)\phi_0(v)
\]
satisfies
\[
 F_m=(N_{m+c_0+1})^d,\qquad \Phi(F_m)\subseteq F_{m+1}.
                                                               \tag{IFD.27}
\]

\begin{proof}
Choose $c_0$ above all pole bounds for the finitely many entries
of $P$; for the empty matrix take $c_0=0$.
Increasing a pole bound preserves it, by multiplying the integral
coefficient by a power of $\tau$.
From $\phi_0(t)\in N_2$ and multiplicativity one gets
$\phi_0(N_r)\subseteq N_{2r}$: apply $\phi_0$ to
$t^rc(b)+pw$, use its preservation of $c(B_0)$, and expand the
$r$th power of the relation in $N_2$.

We also need the following direct product calculation, which does
not require $t$ to be a unit. If $t^{c_0}q=c(b)+pu$ and
$v=t^{r+c_0}c(a)+pw$, then
\[
                         qv=t^rc(ba)+p(t^rc(a)u+qw)\in N_r.
\]
For $v\in F_m$, its Frobenius coordinates are in
$N_{2(m+c_0+1)}\subseteq N_{m+1+2c_0+1}$.
The product calculation lowers the exponent by $c_0$ when
multiplied by any entry of $j_0(P)$. Summing the finitely many
entries gives (IFD.27). In particular $\Phi$ preserves $F_0$.
\end{proof}
% END IFD21

% BEGIN IFD22
\paragraph{IFD22: eventual smallness uses one common iterate.}
\label{claim:IFD22}
Keep IFD21, and assume $\phi_0(j_0(E_0))\subseteq j_0(E_0)$.
Let $S_0\subseteq W_0$ be a subring containing $j_0(E_0)$.
Assume that for every $y\in S_0$ and $m$ there is $y_0$ with
$y-y_0\in N_m$ and $\phi_0^r(y_0)\in j_0(E_0)+pW_0$ for
some $r$. Then for every vector $y\in S_0^d$ there are $r$
and $d_0\in j_0(E_0)^d$ with
\[
                              \Phi^r y-d_0\in F_0.         \tag{IFD.28}
\]

\begin{proof}
Approximate each coordinate at precision $c_0+1$, and take the
maximum $r$ of their finitely many iteration exponents.
Persistence in IFD14 puts every $\phi_0^r((y_0)_i)$ in
$j_0(E_0)+pW_0$. There is a matrix $P_r$ over $E_0$ with
\[
                              \Phi^r(v)=j_0(P_r)\phi_0^r(v).
\]
Indeed $P_0=1$. For the induction, choose preimages in $E_0$ of
the finitely many entries of $\phi_0(j_0(P_r))$, and multiply
the resulting matrix on the left by $P$. This proves the formula
without requiring an inverse to $\phi_0$ or injectivity of $j_0$.
It follows that $\Phi^r(y_0)\in j_0(E_0)^d+pW_0^d$.
Meanwhile $y-y_0\in F_0$, and IFD21 shows
$\Phi^r(y-y_0)\in F_0$. Since $pW_0^d\subseteq F_0$, their
sum proves (IFD.28). Empty coordinate sets give the same statement
with $r=0$.
\end{proof}
% END IFD22

% BEGIN IFD23
\paragraph{IFD23: the full generic matrix quotient theorem.}
\label{claim:IFD23}
Take commutative rings $W_0,E_0,B_0$, maps $j_0,c$, a ring
endomorphism $\phi_0$, a subring $S_0\subseteq W_0$, and
$t=c(\tau)$, with the filtration $N_m$ of IFD1.
Assume precisely the following original matrix hypotheses:
\begin{enumerate}
\item $j_0(E_0)\subseteq S_0$, and $\phi_0$ preserves $S_0$,
      $j_0(E_0)$ and $c(B_0)$.
\item $\phi_0(t)\in N_2$, and every $j_0(e)$ has a bounded
      pole as in IFD21.
\item $S_0$ and $j_0(E_0)$ are $p$-saturated in $W_0$.
\item For every $k$ and sequence $f_i\in S_0\cap N_{i+k}$,
      there is $s\in S_0$ with
      $s-\sum_{i<m}f_i\in N_{m+k}$ for every $m$.
\item An element approximable by $j_0(E_0)$ modulo every
      $N_m$ belongs to $j_0(E_0)+pW_0$.
\item The eventual-coefficient approximation hypothesis of
      IFD22 holds for every element of $S_0$.
\item Every sequence in $S_0$ has a $p$-weighted sum with the
      congruences in property 6 of IFD17; and $j_0(E_0)$ is
      $p$-adically closed in the sense of property 7 there.
\end{enumerate}
For every $d\geq0$ and $P\in M_d(E_0)$, the map
\[
             j_0(P)\phi_0-1:S_0^d/j_0(E_0)^d
                                    \longrightarrow S_0^d/j_0(E_0)^d
                                                               \tag{IFD.29}
\]
is bijective as a map of additive groups. No invertibility of $P$
is required. This is the full original
\texttt{matrix\_frob\_sub\_one\_bijective} implication.

\begin{proof}
Apply IFD20 to the ambient additive group $N=W_0^d$,
$M=S_0^d$, $D=j_0(E_0)^d$, the operator
$\Phi=j_0(P)\phi_0$ and the subgroups $F_m$ in IFD21.
The first hypothesis gives preservation of $M,D$.
IFD21 gives contraction and decreasingness of $F_m$.
Saturation is coordinatewise. Hypothesis 4, with $k=c_0+1$,
supplies the required sums in $M\cap F_m$ one coordinate at a time.
If $x$ is approximable by $D$ modulo every $F_m$, it is
approximable modulo every $N_m$ in each coordinate, because
$N_{m+c_0+1}\subseteq N_m$. Hypothesis 5 gives $x\in D+pN$.
IFD22 gives eventual smallness. Finally the weighted $p$-sums
and closedness in hypothesis 7 apply coordinatewise.
These check every hypothesis of IFD17, so IFD20 proves (IFD.29).
All choices are finite coordinate choices, including the maximum
in IFD22. No extra uniformity over an infinite matrix has entered.
\end{proof}
% END IFD23

% BEGIN IFD24
\paragraph{IFD24: the actual fixed-vector descent and its surjectivity companion.}
\label{claim:IFD24}
For every original good layer and supplied $\Gamma,R,I,\mathrm{IoW}$,
and every $d\geq0$, $P_0\in M_d(E)$, the additive map
\[
              j(P_0)\phi-1:(W^H)^d/j(E)^d
                                      \longrightarrow(W^H)^d/j(E)^d
                                                               \tag{IFD.30}
\]
is bijective. In particular, if $z\in(W^H)^d$ and
$j(P_0)\phi(z)-z\in j(E)^d$, then every coordinate of $z$
belongs to $j(E)$. This includes the full original
\texttt{Rational.vec} theorem.

\begin{proof}
Use IFD23 with the unchanged actual rings and $S_0=W^H$.
IFD7 gives containment of $j(E)$, its Frobenius stability and
the bounded-pole property. WIS20 gives stability of $W^H$,
its $p$-saturation and its weighted $p$-sums.
IFD2 gives stability of $c(B)$ and $\phi(t)\in\mathcal N_2$.
WIS18 gives $p$-saturation of $j(E)$ in $W$: if $px=j(e)$,
reflected divisibility gives $e=pe'$, and absence of $p$-torsion
gives $x=j(e')$.

Since $t=j(a(T))$, IFD7 says $t\in W^H$.
WIS25 therefore supplies the full fixed-ring near-parameter sums
for every shift $k$, using this exact $H$ and this exact $t$.
IFD7 gives near-closedness modulo $p$ and IFD16 gives eventual
coefficient approximation. WIS19 gives $p$-adic closedness of
$j(E)$. These are precisely the seven groups of hypotheses in
IFD23. It proves (IFD.30) and hence the stated vector descent.
The bijectivity assertion is a valid additional consequence of the
same supplied data; the original matrix consumer uses only its
injectivity half.
\end{proof}
% END IFD24

% BEGIN IFD25
\paragraph{IFD25: full inhomogeneous invariant Frobenius-matrix descent.}
\label{claim:IFD25}
Keep every original prime, field, level and supplied record as above.
Let $a,b\geq0$, $A_0\in M_a(E)$, $A_1\in M_b(E)$,
$Y\in M_{a,b}(E)$ and $U\in M_{a,b}(W)$.
Assume $A_0,A_1$ are invertible, every entry of $U$ is $H$-fixed,
and
\[
                        j(A_0)\phi(U)=Uj(A_1)+j(Y).         \tag{IFD.31}
\]
Then every entry of $U$ belongs to $j(E)$.
This is the complete original \texttt{Rational.main}, equivalently
the exported fixed-cyclotomic-kernel Frobenius-solution descent
theorem, including its inhomogeneous term.

\begin{proof}
Choose $C\in M_b(E)$ with $A_1C=1$. Multiplication of (IFD.31)
on the right by $j(C)$ gives
\[
                   j(A_0)\phi(U)j(C)-U=j(YC).
\]
Index the $ab$ entries of $U$ by pairs $(i,j)$ and define
\[
               P_{(i,j),(k,\ell)}=(A_0)_{ik}C_{\ell j}.
\]
If $z_{(i,j)}=U_{ij}$, the preceding equation says exactly
$j(P)\phi(z)-z=j(YC)$ coordinatewise. IFD24 therefore puts
every $z_{(i,j)}$ in $j(E)$, as asserted.
The zero-dimensional cases are vacuous and are included.
The argument in fact only needs a right inverse for $A_1$;
the original invertibility assumption on $A_0$ is retained at
the stated endpoint and has not been traded for a new condition.
\end{proof}
% END IFD25

% BEGIN IFD26
\paragraph{IFD26: exactly one period-ring supplier remains before Wach existence.}
\label{claim:IFD26}
At the original tied local period data, retain the original
crystalline-lattice realization premise used by WLD26.
Suppose the full actual Wach-prism producer
\texttt{exists\_wach\_prism\_over\_integers} is supplied at every
original good layer, including its maps, scalar identity and both
finite coefficient-coordinate assertions.
Then the full original \texttt{WachEx} follows for every permitted
$\Gamma,R,I$ and every supplied original etale module.
The actual Wach prism is the only still-retained period-ring
supplier in this return.

\begin{proof}
WLD26 proves the full original all-layer return from three precisely
listed period-ring suppliers. WIS18 supplies its injectivity and
unit-reflection supplier on the actual $I.\mathrm{ioE}$.
IFD25 with $Y=0$ supplies its homogeneous fixed-matrix descent
for every pair of invertible coefficient matrices and every
rectangular fixed matrix. This is the exact input used in WLD15
and WLD25, with the same coefficient map, Frobenius and kernel.
The remaining prism supplier is the hypothesis printed above.
Thus all three entries of WLD26 are available at each permitted
layer, and its conclusion supplies every field of the original
$\mathrm{NS}$ for every supplied etale module.
No extra compatibility condition between distinct outer layers
is introduced: the original \texttt{WachEx} quantifies over them.

This conditional conclusion does not assert existence of the
retained prism, of the original tied period data, or of the
crystalline-realization theorem. It also does not supply the
separate \texttt{FaithH}, \texttt{EtBasis}, Herr class maps or
quantitative comparison needed for the full $\mathrm{HW/Fon/Cmp}$
package. The common quantitative family in
\texttt{GLOBAL\_DENOMINATOR} remains a later obligation.
\end{proof}
% END IFD26

\paragraph{Where the next proof must begin.}
The coefficient inclusion, both completeness structures and the
entire invariant matrix descent are now available on the original
objects. The next local-ring task is the actual Wach prism with
its prescribed maps and scalar identities. After that task, IFD26
returns the full all-layer lattice existence statement. The
comparison and common-family quantitative arguments remain separate
steps toward the global denominator obligation.
